\documentclass[11pt]{article}

% ============================================================================
%  Modified Forward--Backward objectives on the GCIQL bilinear architecture
%  Living design doc (PAPER/ is gitignored / local-only).
% ============================================================================

\usepackage[margin=1in]{geometry}
\usepackage{amsmath,amssymb,amsthm}
\usepackage{booktabs}
\usepackage{xcolor}
\usepackage{enumitem}
\usepackage[colorlinks=true,linkcolor=blue,citecolor=blue,urlcolor=blue]{hyperref}

% ---- notation ----
\newcommand{\sg}{\operatorname{sg}}                 % stop-gradient
\newcommand{\E}{\mathbb{E}}
\newcommand{\R}{\mathbb{R}}
\newcommand{\Ind}{\mathbf{1}}
\newcommand{\D}{\mathcal{D}}
\newcommand{\Sspace}{\mathcal{S}}
\newcommand{\Aspace}{\mathcal{A}}
\newcommand{\inner}[2]{\langle #1,\, #2\rangle}
\newcommand{\Fmap}{F}                                % forward map  = critic phi
\newcommand{\Bmap}{B}                                % backward map = critic psi
\newcommand{\Lexp}{L^{\tau}_2}                       % expectile loss
\newcommand{\good}[1]{\textcolor{teal}{#1}}
\newcommand{\bad}[1]{\textcolor{red!70!black}{#1}}

\title{\textbf{Modified Forward--Backward Objectives on the\\
GCIQL Bilinear Architecture}\\[4pt]
\large What the bridge experiments show, and how to put policy improvement
back into an FB-like objective}
\author{Factored-FB / \texttt{gciql-fb-bridge}}
\date{2026-07-01}

\begin{document}
\maketitle

\begin{abstract}
Our FB${}\leftrightarrow{}$GCIQL bridge experiments morph a working
goal-conditioned IQL agent (GCIQL, bilinear critic, ${\sim}90\%$ on
cube-single) toward Forward--Backward (FB) along a set of \emph{design-decision
clusters} that separate FB's \emph{architecture} from FB's \emph{objective}.
The result is sharp: FB's full architectural skeleton survives on the strong
IQL base (cluster B endpoint \good{$87\%$}), but swapping the IQL value backup
for FB's successor-measure objective on the \emph{same} architecture collapses
performance (cluster C endpoint \bad{$16\%$}). The single distinguishing
ingredient is a \emph{policy-improvement operator}: IQL's expectile backup
performs in-sample improvement; FB's measure loss is an on-policy
\emph{evaluation} objective with no maximization. This note (i) writes down what
we observe, (ii) diagnoses the collapse, and (iii) proposes two directions --- a
zero-shot multi-task \emph{in-sample Forward--Backward} objective and a
single-task occupancy-control ladder --- laying out the loss modifications
mathematically for each.
\end{abstract}

% ============================================================================
\section{Setup and notation}
% ============================================================================

Offline dataset $\D=\{(s,a,s')\}$ over states $s\in\Sspace$, actions
$a\in\Aspace$, discount $\gamma$. Goals $g$ and task vectors $z\in\R^{d}$ are
sampled by hindsight relabeling from future states; $\rho$ denotes the base
state distribution. Latent dimension $d$ (we use $d{=}50$).

\paragraph{GCIQL bilinear critic (the ${\sim}90$ base).}
The critic factorizes the goal-conditioned action value as a normalized inner
product of a \emph{forward} map $\Fmap$ and a \emph{backward} map $\Bmap$:
\begin{equation}
  Q(s,a,g) \;=\; \tfrac{1}{\sqrt d}\,\inner{\Fmap(s,a)}{\Bmap(g)},
  \qquad \Fmap,\Bmap:\;\to\R^{d}.
  \label{eq:bilinear-q}
\end{equation}
(In code $\Fmap=\phi$, $\Bmap=\psi$; two critic heads give the double-$Q$.) The
value $V(s,g)$ is an MLP (or bilinear). A one-step \emph{flow} actor $\pi$ is
distilled from a multi-step flow policy and trained to maximize $Q$.

\paragraph{IQL objective (cluster B).}
With sparse goal reward $r(s,g)=\Ind[s=g]-1\in\{-1,0\}$ (\texttt{gc\_negative})
and termination mask $m$, GCIQL trains
\begin{align}
  \mathcal{L}_{V}
    &= \E_{(s,a,g)}\!\big[\Lexp\!\big(\bar Q(s,a,g)-V(s,g)\big)\big],
    &\Lexp(u) &= \big|\tau-\Ind[u<0]\big|\,u^{2},
    \label{eq:iql-v}\\[2pt]
  \mathcal{L}_{Q}
    &= \E_{(s,a,s',g)}\!\big[\big(r(s,g)+\gamma\,m\,V(s',g)-Q(s,a,g)\big)^{2}\big],
    \label{eq:iql-q}
\end{align}
where $\bar Q$ uses the target critic. The \emph{expectile} $\tau=0.9$ makes $V$
approximate $\max_{a\in\mathrm{supp}(\pi_\beta)}Q(s,a,g)$: an \good{in-sample
policy-improvement operator} that never queries out-of-distribution actions.
The flow actor then goes greedy on the improved $Q$.

\paragraph{Forward--Backward objective (cluster C).}
FB parameterizes the \emph{successor measure} of policy $\pi_z$ as the same
inner product,
\begin{equation}
  M^{\pi_z}(s,a,s') \;=\; \tfrac{1}{\sqrt d}\,\inner{\Fmap(s,a,z)}{\Bmap(s')}
  \;\approx\; \sum_{t\ge 0}\gamma^{t}\,\Pr\!\big(s_{t+1}=s'\mid s,a,\pi_z\big)
  \;\big/\;\rho(s'),
  \label{eq:succ-measure}
\end{equation}
i.e.\ the discounted future-state occupancy density w.r.t.\ $\rho$. Any reward
reads out linearly through a task encoding, giving FB its zero-shot property:
\begin{equation}
  z_r=\E_{s\sim\rho}\!\big[r(s)\,\Bmap(s)\big],
  \qquad
  Q_{z}(s,a)=\inner{\Fmap(s,a,z)}{z},
  \qquad\text{(goal-reaching: } z=\Bmap(g)\text{)}.
  \label{eq:fb-readout}
\end{equation}
On a minibatch $\{(s_i,a_i,s_i')\}_{i=1}^{B}$ with $z_i=\sg\,\Bmap(g_i)$, let
$M_{ij}=\tfrac1{\sqrt d}\inner{\Fmap(s_i,a_i,z_i)}{\Bmap(s_j')}$ and the
\emph{target} $\tilde M_{ij}=\tfrac1{\sqrt d}\inner{\Fmap(s_i',a_i^{+},z_i)}{\Bmap(s_j')}$
with a bootstrap action $a_i^{+}\sim\pi(\cdot\mid s_i',z_i)$ from the target net
(stop-grad). Writing $\Delta_{ij}=M_{ij}-\gamma\tilde M_{ij}$ and
$\mathrm{off}=\{i\ne j\}$:
\begin{align}
  \mathcal{L}_{\mathrm{off}}
    &= \frac{1}{2\,|\mathrm{off}|}\sum_{i\ne j}\Delta_{ij}^{2},
    &\text{(off-diagonal measure TD)}
    \label{eq:fb-off}\\
  \mathcal{L}_{\mathrm{diag}}
    &= -\frac{1}{B}\sum_{i}\Delta_{ii},
    &\text{(diagonal: observed successor is reachable)}
    \label{eq:fb-diag}\\
  \mathcal{L}_{\perp}
    &= \frac{1}{2\,|\mathrm{off}|}\sum_{i\ne j}\inner{\Bmap(s_i')}{\Bmap(s_j')}^{2}
       -\frac1B\sum_i\lVert\Bmap(s_i')\rVert^{2},
    &\text{(orthonormality } \lVert\E\Bmap\Bmap^{\!\top}-I\rVert)
    \label{eq:fb-ortho}
\end{align}
and $\mathcal{L}_{\mathrm{FB}}=\mathcal{L}_{\mathrm{off}}+\mathcal{L}_{\mathrm{diag}}+\lambda\mathcal{L}_{\perp}$.
The greedy flow actor maximizes $Q_z(s,a)=\inner{\Fmap(s,a,z)}{z}$.

\medskip
\noindent\textbf{The structural point.} Equations~\eqref{eq:bilinear-q} and
\eqref{eq:succ-measure} are the \emph{same} bilinear form. Clusters B and C hold
that architecture fixed and vary only the objective
(\ref{eq:iql-v}--\ref{eq:iql-q}) vs.\ (\ref{eq:fb-off}--\ref{eq:fb-ortho}).

% ============================================================================
\section{What the bridge experiments show}
% ============================================================================

We morph GCIQL toward FB along ten design axes grouped into clusters. The two
that matter here are orthogonal:
\begin{itemize}[leftmargin=1.4em,itemsep=2pt]
  \item \textbf{Cluster B --- ``FB skeleton, IQL blood.''} Bolt FB's
        \emph{architecture} onto the winning IQL critic, keeping the IQL backup:
        two-tower forward map $\Fmap=\mathrm{trunk}(\mathrm{MLP}[s,a],\mathrm{MLP}[s,z])$,
        hard sphere normalization $\lVert z\rVert=\lVert\Bmap\rVert=\sqrt d$
        (\texttt{project\_z}), and FB orthonormality at its true coefficient
        $\lambda=1$. Endpoint \textbf{B4} $=$ the full stack.
  \item \textbf{Cluster C --- ``rebuild FB inside GCIQL.''} Hold the architecture
        and swap the objective for FB's successor measure
        (\ref{eq:fb-off}--\ref{eq:fb-ortho}), adding restrictions one rung at a
        time. Endpoint \textbf{C4} $=$ measure $+$ $z$-conditioned forward map
        $+$ sphere-norm $+$ $\lambda=1$.
  \item \textbf{Cluster D --- conditioning source.} Route the actor's task
        vector $z$ from the \emph{critic} $\Bmap$ (not the value net). Endpoint
        \textbf{D1}.
\end{itemize}

\begin{table}[h]
\centering
\small
\begin{tabular}{@{}llll r@{}}
\toprule
\textbf{Config} & \textbf{Objective} & \textbf{Architecture} & \textbf{Note} & \textbf{Success} \\
\midrule
GCIQL-Bilinear-$d50$          & IQL       & bilinear                 & the ${\sim}90$ base            & \good{$90.7\pm2$} \\
\;+ two-tower (alone)         & IQL       & \;+ two-tower            & FB fwd map, nearly free        & $86.4\pm2$ \\
GCIQL-MLP-value               & IQL       & MLP critic               & no bilinear factorization      & $73$ \\
\midrule
\textbf{B4} FB-struct         & IQL       & two-tower FB skeleton    & two-tower+projZ+ortho$_1$      & \good{$\mathbf{87.1}$} \\
\;\;B-TwoTower-ProjZ (abl.)   & IQL       & two-tower+projZ          &                                & $87.5$ \\
\;\;B-TwoTower-Ortho$_1$ (abl.)& IQL      & two-tower+ortho$_1$      &                                & $85.5$ \\
\midrule
\textbf{C4} FB-rebuild        & measure   & $z$-fwd FB skeleton      & the objective swap             & \bad{$\mathbf{16.5}$} \\
\;\;C2-MeasureTwoTower        & measure   & two-tower FB skeleton    & \emph{architecture-matched to B4} & \bad{${\sim}8.0$} \\
\;\;C-FBz-O$_1$ (abl.)        & measure   & \;+ $z$-fwd, ortho$_1$   &                                & $17.1$ \\
\;\;C-FBz-ProjZ-O$_{1000}$    & measure   & \;+ projZ, ortho$_{1000}$&                                & $18.0$ \\
GCIQL-Bil-FBz-O$_{1000}$      & measure   & bilinear + $z$-fwd       & measure + genuine $F(s,a,z)$   & $25.5$ \\
GCIQL-Bil-MeasureOnly         & measure   & bilinear                 & bare measure graft             & \bad{$9.9$} \\
\midrule
\textbf{D1} actor-$z$-from-$\Bmap$ & IQL$^\dagger$ & bilinear            & conditioning-source fix        & $22.0$ \\
plain FB (reference, torch)   & measure   & native FB                & the FB anchor                  & $48.0$ \\
\bottomrule
\end{tabular}
\caption{Cube-single-state, $\mathrm{evaluation/overall\_success}$ @500k
($n{=}3$ unless noted). $^\dagger$D1 pairs the IQL/measure cells with a
critic-sourced actor $z$. FB-side morphs (start from FB, import IQL's value
objective; not shown) span $48$--$65$ and \emph{never} reach $90$.}
\label{tab:bridge}
\end{table}

\paragraph{Reading the table.}
\begin{enumerate}[leftmargin=1.5em,itemsep=2pt]
  \item \textbf{Architecture is nearly free.} Every cluster-B cell sits at
        $85$--$88$, within a few points of the unmodified $90.7$. The full FB
        skeleton (two-tower $+$ sphere-norm $+$ orthonormality) survives on the
        strong base. The only real cost is two-tower alone ($86.4$ vs $90.7$);
        the remaining knobs add essentially nothing.
  \item \textbf{The objective is the killer.} The \emph{same} architecture under
        the measure objective collapses: C4 $=16.5$, and no measure ablation
        (ortho $1{\to}1000$, $\pm$sphere-norm, $\pm z$-forward) escapes the
        $9$--$25$ band. Adding a genuine $z$-indexed forward family lifts the
        floor only to $25.5$.
  \item \textbf{Neither direction reaches $90$ from the FB side.} Starting from
        native FB and importing IQL restrictions plateaus at $48$--$65$. The
        lever is the base/objective, not FB's structure --- confirmed from both
        ends.
\end{enumerate}

% ============================================================================
\section{Diagnosis: the missing policy-improvement operator}
% ============================================================================

The measure loss (\ref{eq:fb-off}--\ref{eq:fb-ortho}) is a \emph{policy
evaluation} objective. Its bootstrap $a_i^{+}\sim\pi(\cdot\mid s_i',z_i)$ is
on-policy (SARSA-like), and the diagonal term \eqref{eq:fb-diag} merely asserts
``the observed successor $s_i'$ is reachable from $(s_i,a_i)$'' ---
reward-agnostic. \emph{There is no maximization anywhere.} The object it fits is
the successor measure of the current/behavior policy; the greedy flow actor over
$Q_z=\inner{\Fmap}{z}$ then extracts a policy no better than the data supports.

IQL \eqref{eq:iql-v}, by contrast, decouples an \emph{improved} value $V\approx
\max_a Q$ via the expectile and regresses $Q$ toward $r+\gamma V$. That single
in-sample-max is what B4 keeps and C4 discards.

\begin{center}\itshape
Same bilinear factorization, same FB skeleton: $87\%$ with an improvement
operator (B4), $16\%$ without one (C4). The design target is therefore an
FB-like \emph{occupancy} objective that carries an in-sample improvement
operator.
\end{center}

This is consistent with, and narrower than, the paper's representation thesis
(occupancy under-resolves the interaction variable): B4$=87$ proves the
\emph{architecture} has the capacity, so any residual shortfall is
objective/optimization, not a representational ceiling.

% ============================================================================
\section{The core modification: in-sample improvement in the measure backup}
% ============================================================================

IQL's trick is to introduce a \emph{state} value $V(s)$ fit as an expectile of
the \emph{state--action} value $Q(s,a)$ over dataset actions, then bootstrap
through $V$ (never through a sampled $a^{+}$). We port this \emph{to the measure
level}. Introduce a \textbf{value successor measure}
\begin{equation}
  N^{\pi}(s,s') \;=\; \tfrac{1}{\sqrt d}\,\inner{\Fmap_V(s,z)}{\Bmap(s')},
  \label{eq:value-measure}
\end{equation}
a state-only occupancy with its own forward head $\Fmap_V$ (the ``$V$-leg'',
already present as \texttt{measure\_legs='value'}). We want $N$ to be the
occupancy of the \emph{expectile-improved} policy. The key observation that
makes a \emph{vector} expectile well defined: each backward direction
$\Bmap(s_j')$ in the batch defines a scalar occupancy
$\inner{\Fmap(s,a,z)}{\Bmap(s_j')}$, so applying the scalar expectile
\emph{per column $j$} imposes it along $B$ sampled directions
simultaneously --- the off-diagonal batch structure of
(\ref{eq:fb-off}) already \emph{is} the sampling over directions $z$.

Concretely, with $M_{ij}=\tfrac1{\sqrt d}\inner{\Fmap(s_i,a_i,z_i)}{\Bmap(s_j')}$
(action measure) and $N_{ij}=\tfrac1{\sqrt d}\inner{\Fmap_V(s_i,z_i)}{\Bmap(s_j')}$
(value measure):
\begin{align}
  \mathcal{L}_{N}
    &= \frac{1}{B^2}\sum_{i,j}\Lexp\!\big(\,\sg M_{ij}-N_{ij}\,\big),
    &\text{(expectile: } N \to \text{in-sample }\max_a M\text{ per direction)}
    \label{eq:isfb-v}\\
  \mathcal{L}_{M}
    &= \frac{1}{2\,|\mathrm{off}|}\sum_{i\ne j}\!\big(M_{ij}-\gamma\,\sg N^{-}_{ij}\big)^{2}
       \;-\;\frac1B\sum_i\big(M_{ii}-\gamma\,\sg N^{-}_{ii}\big),
    &\text{(measure TD through } N,\text{ no } a^{+})
    \label{eq:isfb-m}
\end{align}
where $N^{-}$ uses the target $\Fmap_V$. The orthonormality term
\eqref{eq:fb-ortho} is retained. This is exactly IQL's $(Q,V)$ decoupling
(\ref{eq:iql-v}--\ref{eq:iql-q}) with scalars replaced by occupancy inner
products, and it \good{eliminates the on-policy bootstrap} entirely: the target
in \eqref{eq:isfb-m} is the improved value measure, never a sampled OOD action.
We call this \textbf{in-sample Forward--Backward (IS-FB)}. Everything below is a
specialization or alternative of (\ref{eq:isfb-v}--\ref{eq:isfb-m}).

% ============================================================================
\section{Direction 1 --- zero-shot multi-task, reach ${\sim}90$}
% ============================================================================

\emph{Goal:} preserve FB's actual prize --- one network, task $z_r=\E_\rho[r\Bmap]$
inferred at test time for \emph{any} reward with no fine-tuning
\eqref{eq:fb-readout} --- while reaching IQL-level control.

\subsection*{1A. In-sample Forward--Backward (IS-FB)}
Train (\ref{eq:isfb-v}--\ref{eq:isfb-m})$+\lambda\mathcal{L}_{\perp}$ with the
$z$-conditioned two-tower forward map (the B4 architecture, now under the measure
objective). Directions $\Bmap(s_j')$ are drawn from \emph{full-support} future
states (not only the current goal), so the fitted measure covers arbitrary $z$;
at test time any reward reads out via \eqref{eq:fb-readout}. This is the minimal
change that gives the measure a policy-improvement operator while keeping the
zero-shot factorization intact.
\begin{itemize}[leftmargin=1.4em,itemsep=1pt]
  \item \emph{Discriminates:} does injecting IQL's in-sample-max into the
        successor measure recover the C4$\to$B4 gap ($16\to{\sim}90$) \emph{and}
        retain zero-shot transfer?
  \item \emph{Risk:} multi-task fitting dilutes improvement across the whole
        $z$-family; the expectile is imposed per-direction but the shared
        $\Fmap_V$ must serve all of them. (This is precisely the confound that
        Direction~2 removes.)
\end{itemize}

\subsection*{1B. Successor Features $+$ GPI (the ``SF not SM'' pivot)}
FB fits occupancy of \emph{raw states} ($\Bmap=\Bmap(s')$). Successor Features
(SF) instead fit occupancy of \emph{features} $\varphi(s)$, learn task weights
$w$ by reward regression, and improve by generalized policy improvement:
\begin{align}
  \psi^{\pi}(s,a) &= \E\!\Big[\textstyle\sum_{t\ge0}\gamma^{t}\varphi(s_{t})\;\Big|\;s_0{=}s,a_0{=}a,\pi\Big],
   &\text{(TD: } \psi(s,a)=\varphi(s)+\gamma\,\sg\,\psi(s',a'),\; a'\sim\pi) \label{eq:sf}\\
  r(s) &\approx \inner{\varphi(s)}{w},
   &\text{(reward regression; } Q^{\pi}_w=\inner{\psi^\pi}{w}) \label{eq:sf-reward}\\
  \pi_{\mathrm{GPI}}(s) &= \arg\max_{a}\;\max_{k}\;\inner{\psi^{\pi_k}(s,a)}{w}.
   &\text{(generalized policy improvement)} \label{eq:sf-gpi}
\end{align}
SF$+$GPI keeps the bilinear occupancy factorization and zero-shot linear-reward
readout, but adds a \good{genuine reward signal in the loop}
\eqref{eq:sf-reward} and a \good{real max operator} \eqref{eq:sf-gpi}. On the B4
architecture: $\psi=$ two-tower forward map, $\varphi=$ base features (a frozen
or slowly-learned $\Bmap$), $w=z$.
\begin{itemize}[leftmargin=1.4em,itemsep=1pt]
  \item \emph{Discriminates a different hypothesis:} is FB's failure that
        \emph{raw-state} occupancy is too hard to fit on manipulation, versus
        \emph{feature} occupancy with reward-in-loop?
  \item \emph{Trade-off:} GPI needs a set of base policies $\{\pi_k\}$; with a
        single policy it degenerates to on-policy SF (same evaluation trap as
        FB). The improvement therefore comes from the policy set, which must be
        built (e.g.\ per-$z$ greedy policies).
\end{itemize}

% ============================================================================
\section{Direction 2 --- single-task occupancy control (isolate the operator)}
% ============================================================================

\emph{Goal:} drop multi-task. Pin $z=\Bmap(g)$ to the goal, and ask the clean
question the C4-vs-B4 result poses: \emph{can an occupancy-fitting objective plus
an improvement operator drive the B4 representation to $90$, with the
$z$-family-dilution confound removed?} A cheap$\to$expensive ladder, each rung
reusing the last.

\subsection*{2B. Reward-shaped measure (cheapest probe)}
The diagonal term \eqref{eq:fb-diag} attracts occupancy at the observed
successor \emph{regardless of reward}. Gate it by the task reward so that only
\emph{reward-relevant} occupancy is pulled up. With per-sample weight
$w_i=\exp\!\big(\beta\,A(s_i,a_i,g_i)\big)$ (an advantage/expectile weight, as in
AWR) or simply $w_i=\Ind[\text{$s_i'$ advances toward $g_i$}]$:
\begin{equation}
  \tilde{\mathcal{L}}_{\mathrm{diag}}
    = -\frac{1}{\sum_i w_i}\sum_{i} w_i\,\big(M_{ii}-\gamma\,\sg\tilde M_{ii}\big),
  \qquad A(s,a,g)=Q(s,a,g)-V(s,g).
  \label{eq:rew-shaped}
\end{equation}
This is the minimal, one-line change to \texttt{measure\_loss}: soft policy
improvement by reweighting which transitions' occupancy is reinforced.
\emph{Discriminates:} how much of the gap closes from reward-shaping alone,
before any structural expectile.

\subsection*{2C. Value-legs expectile (small delta, reuses existing path)}
The state-only measure ($\texttt{measure\_legs='value'}$) already fits
$N(s,s')=\tfrac1{\sqrt d}\inner{\Fmap_V(s)}{\Bmap(s')}$ against
$\gamma N(s',s')$. Replace its symmetric TD with the \emph{expectile} backup so
$N$ tracks the occupancy of an improved policy:
\begin{equation}
  \mathcal{L}_{N}^{\,\mathrm{2C}}
    = \frac{1}{2\,|\mathrm{off}|}\sum_{i\ne j}
      \Lexp\!\big(\,M_{ij}-\gamma\,\sg N^{-}_{ij}\,\big)
    \;-\;\frac1B\sum_i\big(M_{ii}-\gamma\,\sg N^{-}_{ii}\big),
  \label{eq:vleg-expectile}
\end{equation}
i.e.\ apply $\Lexp$ (rather than plain $\ell_2$) to the off-diagonal measure
residual, biasing the fit toward the upper occupancy envelope. Cheapest code
delta that introduces a max-like operator.

\subsection*{2A. In-sample FB, single goal (the clean test)}
Equations~(\ref{eq:isfb-v}--\ref{eq:isfb-m}) with $z=\Bmap(g)$ fixed. Because
directions in the batch are exactly the goal-reaching family, this is precisely
IQL \eqref{eq:iql-v}--\eqref{eq:iql-q} with the MLP $Q,V$ replaced by the
occupancy inner products $M,N$ --- ``IQL whose value is a successor measure.''
\begin{itemize}[leftmargin=1.4em,itemsep=1pt]
  \item \emph{Discriminates the diagnosis cleanly:} if 2A reaches ${\sim}90$, the
        occupancy objective \emph{per se} was never the problem --- the missing
        improvement operator was, and Direction~1 need only scale the same recipe
        to full-support $z$. If 2A plateaus below B4, the occupancy objective is
        fundamentally weaker than reward-TD even with zero dilution, and the SF
        pivot (1B) or a hybrid becomes the fallback.
  \item \emph{Relation to 1A:} 2A \emph{is} 1A restricted to the goal-reaching
        reward family; the training objective is identical, only the direction
        sampling (goal vs.\ full-support) and the evaluation (single goal vs.\
        arbitrary $z_r$) differ. Hence a 2A win makes 1A nearly free.
\end{itemize}

\paragraph{Relation to the bridge's D8 (goal/task distribution).} The
direction-sampling choice above (goal-pinned in 2A vs.\ full-support in 1A) \emph{is}
the bridge's D8 axis --- the distribution of $z$ the objective is trained on. It is
already folded into the 1A/2A split, so no separate knob is needed. The other half of
D8, \emph{decoupling the actor's task distribution from the critic's}, is by contrast
a \B{closed negative} on the bridge (zHER $32$, decouple-actor-$z$ $18$--$39$, all
below plain FB $48$): it does not rescue a $\max$-free measure and is not proposed
here. In-sample FB adds the $\max$ operator D8 never touched; the two axes are
orthogonal.

% ============================================================================
\section{Sequencing and expected outcomes}
% ============================================================================

\begin{table}[h]
\centering
\small
\begin{tabular}{@{}clllc@{}}
\toprule
\textbf{Order} & \textbf{Experiment} & \textbf{Modification} & \textbf{Discriminates} & \textbf{Target} \\
\midrule
1 & 2B reward-shaped     & \eqref{eq:rew-shaped}                       & reward-shaping alone        & $\uparrow$ from $16$ \\
2 & 2C value-leg exp.\   & \eqref{eq:vleg-expectile}                   & cheap max operator          & $\uparrow$ \\
3 & 2A IS-FB (single)    & (\ref{eq:isfb-v}--\ref{eq:isfb-m}), $z{=}\Bmap(g)$ & \textbf{operator hypothesis} & ${\sim}90$ (vs B4 $87$) \\
4 & 1A IS-FB (multi)     & (\ref{eq:isfb-v}--\ref{eq:isfb-m}), full-support $z$ & zero-shot + control   & ${\sim}90$ zero-shot \\
5 & 1B SF$+$GPI          & (\ref{eq:sf}--\ref{eq:sf-gpi})              & raw-state vs feature occ.\  & (parallel alt.) \\
\bottomrule
\end{tabular}
\caption{Recommended order. Do Direction~2 first \emph{even though the prize is
Direction~1}: 1A's objective equals 2A's plus multi-task direction sampling, so
running 1A first would conflate ``missing improvement'' with ``multi-task
dilution.'' 2A tests the operator hypothesis with one fewer confound and hands 1A
a proven recipe. 1B is a parallel fallback if the SM ladder stalls.}
\label{tab:sequence}
\end{table}

\paragraph{Falsifiable predictions.}
\begin{itemize}[leftmargin=1.4em,itemsep=2pt]
  \item If the diagnosis is right, \textbf{2A reaches ${\sim}85$--$90$} (matching
        B4): occupancy $+$ in-sample-max $\equiv$ IQL. This would be the decisive
        result --- the FB objective was missing \emph{only} an improvement
        operator, not anything representational.
  \item \textbf{2A plateauing at $16$--$50$} would instead implicate the
        occupancy \emph{representation} on manipulation (the paper thesis): even
        with improvement, fitting raw-state successor measure under-resolves the
        interaction variable, and SF (1B) or reward-TD is required.
  \item \textbf{1A tracking 2A} (${\sim}90$) would deliver the paper's
        constructive counterpart: a single zero-shot FB-style network that
        matches goal-conditioned IQL --- ``FB fails on manipulation \emph{because
        of the objective, and here is the objective that fixes it.}''
\end{itemize}

% ============================================================================
\section*{Appendix: implementation touch-points}
% ============================================================================
\small
All modifications live in \texttt{tools/wandb\_mode\_shim/gciql\_flow.py}:
\begin{itemize}[leftmargin=1.4em,itemsep=1pt]
  \item \textbf{IS-FB (1A/2A):} add a value-leg forward head $\Fmap_V$ (or reuse
        the \texttt{value} net's $\phi$), implement $\mathcal{L}_N$
        \eqref{eq:isfb-v} and the $N$-bootstrapped $\mathcal{L}_M$
        \eqref{eq:isfb-m} in \texttt{measure\_loss}; drop the $a^{+}\sim\pi$
        target. New knob \texttt{measure\_insample=True}, expectile reuses
        \texttt{config.expectile}.
  \item \textbf{2B:} weight the diagonal by an advantage/expectile weight
        \eqref{eq:rew-shaped}; needs $V(s,g)$ available under \texttt{measure\_only}.
  \item \textbf{2C:} swap $\ell_2$ for \texttt{expectile\_loss} on the value-leg
        off-diagonal residual \eqref{eq:vleg-expectile}.
  \item \textbf{1B (SF):} new forward-map TD toward $\varphi(s)+\gamma\psi(s',a')$
        \eqref{eq:sf} $+$ reward-regression head for $w$ $+$ GPI over a small
        policy set at action selection.
\end{itemize}
TDD each loss (finite gradients, $Q$/measure rises without critic inflation ---
cf.\ the frozen-critic rule in the actor $Q$-term) before any launch.

\end{document}
